\section{Performance Evaluation}
\subsection{Evidence audit}
Table~\ref{tab:evidence} distinguishes source implementation, tests, and
committed experimental evidence. This distinction is essential because a
policy constructor or test does not establish a performance improvement.

\begin{table}[t]
\centering
\caption{Evidence status at the audited repository revision.}
\label{tab:evidence}
\scriptsize
\resizebox{\columnwidth}{!}{%
\begin{tabular}{p{.31\columnwidth}p{.18\columnwidth}p{.39\columnwidth}}
\toprule
Capability & Status & Evidence boundary\\
\midrule
Multi-device topology & implemented/tested & explicit topology code and tests; no final campaign\\
Multi-edge/shared cloud & implemented/tested & backend and topology configuration; no final campaign\\
Random, greedy, EDF & implemented/validated & committed three-seed synthetic CSV; controlled scenario tests\\
Device/edge/cloud RL & implemented/tested & adapters, observations, credit tests; no committed run\\
Centralized RL & implemented/tested & centralized adapter and tests; no committed run\\
Hierarchical coordination & implemented/tested & causal feedback unit tests; no final ablation\\
Alibaba workload & adapter/tested & mapping and loader tests; no committed trace results\\
Feedback overhead & modeled/tested & counters and theoretical delay; non-blocking\\
\bottomrule
\end{tabular}
}
\end{table}

\subsection{Committed baseline results}
The same seeded workload is passed to every policy. Table~\ref{tab:baselines}
reports the mean across three seeds and the sample standard deviation from
\texttt{results/dev-smoke/aggregate\_statistics.csv}.

\begin{table}[t]
\centering
\caption{Three-seed synthetic baseline evidence (mean $\pm$ standard deviation).}
\label{tab:baselines}
\scriptsize
\resizebox{\columnwidth}{!}{%
\begin{tabular}{lrrrr}
\toprule
Policy & Mean L (s) & P95 L (s) & Deadline & Energy (J)\\
\midrule
Local & 3.164 $\pm$ .467 & 5.491 $\pm$ .612 & 36.7\% & .619 $\pm$ .072\\
Greedy/edge & .602 $\pm$ .118 & .927 $\pm$ .137 & 100\% & .565 $\pm$ .165\\
Edge execute & .602 $\pm$ .118 & .927 $\pm$ .137 & 100\% & .565 $\pm$ .165\\
Cloud/EDF & .513 $\pm$ .113 & .745 $\pm$ .160 & 100\% & .565 $\pm$ .165\\
Cloud/FIFO & .513 $\pm$ .113 & .745 $\pm$ .160 & 100\% & .565 $\pm$ .165\\
Random & 1.177 $\pm$ .387 & 2.625 $\pm$ .761 & 90.0\% & .424 $\pm$ .067\\
\bottomrule
\end{tabular}
}
\end{table}

Cloud forwarding has the lowest mean latency in this light synthetic
workload, but it doubles the modeled payload traffic relative to edge
execution (27.48 MB versus 13.74 MB per seed-average aggregate). Local
execution consumes more device CPU energy and misses more deadlines. These
observations are scenario-specific: greedy/edge equals edge execution because
all default tasks satisfy the greedy offload inequality, and EDF equals FIFO
because the smoke workload does not create a meaningful cloud ordering choice.
The controlled contention scenario distinguishes their queue order but still
misses its deliberately urgent task under both policies.

\subsection{Unvalidated comparisons}
The repository contains scripts for centralized RL, hierarchical RL, three
feedback modes, and synthetic load/network scenarios. It does not contain the
models, comparison CSVs, or run logs needed to reproduce those measurements at
the audited revision. Accordingly, no RL ranking, convergence, topology
sensitivity, feedback benefit, or Alibaba improvement is reported. The
required next experiment is to run training and held-out evaluation with
disjoint seeds, then export the configured metrics and metadata alongside the
commit and package versions.
